% ---
% id: meat-core/rem:witness-exactness
% kind: rem
% title: "Exactness of the exhibited weights and states"
% summary: >-
%   A witness is computed, not declared, so its weights and states can be made
%   exact by construction; a naively summed weight vector can read 1 while
%   its real sum exceeds 1. The diagonal form of (W4) is one exact
%   presentation of a state, not the only one binary64 admits.
% prominence: standard
% status: accepted
% uses: [def:classes, thm:removal]
% mentions: [def:declaration]
% assumes: []
% refs: []
% provenance: newly-derived
% lean:
% ---

Conditions (W3) and (W4) of Definition~\ref{def:classes} read exhibited
\tref{binary64 numbers}, and the normalization step of Theorem~\ref{thm:removal}
uses $\sum_i\mu_i=1$ and $\tr\chi_i=1$ as identities of reals.

\begin{remark}\label{rem:witness-exactness}
\emph{Exact by construction.} (W3) and (W4) are decided at zero defect,
but not in the way the declared marginal $\tau_A$ is checked at admission.
That check decides \tref{declared data}, which an \tref{instance}
\tref[snap]{snaps} under a \tref{producer}'s obligation and a
\tref{consumer} must take as it finds them. A witness is instead a
\tref{computed object} that the implementation exhibits, and the
implementation can make both conditions exact by construction --- which
(W3) and (W4) as stated require it to. The difference is on the false-positive side, which is the
unsafe one. A naively summed $\mu=(1,2^{-60})$ returns the datum $1.0$
while $\sum_i\mu_i>1$ as reals, and then the normalization step of
Theorem~\ref{thm:removal} reads
$\tr\hat\omega=\tr\omega+(\sum_i\mu_i-1)\tr m\ne1$: the modified point is
not feasible and the removal would have been licensed on an identity that
does not hold. Requiring the sum in the exhibited order and verified to
round nowhere closes that window exactly as (W6) closes it for the
residuals, and requiring the exhibited form of (W4) closes the
corresponding window on a computed $\lmin(\chi_i)$ or a computed
$\tr\chi_i$. \emph{Neither leaves a residual exposure}: on the forms those
two conditions admit, both hold of the real numbers the data denote.

\emph{The diagonal form is one exact presentation among others.} Any
presentation deciding $\chi_i\succeq0$ and $\tr\chi_i=1$ exactly by
construction would serve, and the diagonal form is \emph{one} such
presentation, not the only one binary64 admits. An exhibited unit vector's
outer product $\dyad{\psi_i}$ with $\braket{\psi_i}{\psi_i}$ verified to
round nowhere is a genuine second: on a four-dimensional $\Qsys_{s_i}$ the
vector $\psi_i=(\tfrac12,\tfrac12,\tfrac12,\tfrac12)$ has exactly
representable entries, its norm accumulates $0.25\to0.5\to0.75\to1.0$
rounding nowhere, and its outer product --- every entry $\tfrac14$ --- is a
rank-one state that is \emph{not} diagonal. (W4) does not admit it, and
excluding an admissible presentation can only \emph{refuse} a removal. No
impossibility is claimed here.

\emph{The evaluation order of (W6).} The order in which a witness's
residuals are formed is exhibited with the witness. This is the one
difference from the \tref{declared order} of \eqref{eq:declorder}, which is
part of the declaration because the datum it governs is declared.
\end{remark}
